\chapter{Suites, régularités et aptitudes}\label{ch:iv1:sequences-patterns-and-aptitude}

\begin{ivtrains}
Les questions d'un premier tour de test en ligne~: le terme suivant d'une suite, une régularité énoncée en mots, un court raisonnement logique et un calcul rapide à partir d'un petit tableau. Chaque question repose sur une seule règle, plus simple que toute règle concurrente, et une part de la compétence consiste à savoir quand on l'a trouvée. Il faut se chronométrer~: une épreuve réelle laisse bien en deçà du temps prévu par les étoiles.
\end{ivtrains}

\begin{question}[$\star$ \iqroles{QT, Trader, Analyste} \iqfirm{teneur de marché, banque} \iqlevel{junior}]\label{iq:iv1:sequences-patterns-and-aptitude:1}
Quel est le nombre suivant dans la suite 3, 6, 18, 72~?
\end{question}

\begin{question}[$\star$ \iqroles{QT, Trader, Analyste} \iqfirm{teneur de marché, banque} \iqlevel{junior}]\label{iq:iv1:sequences-patterns-and-aptitude:2}
Quel est le nombre suivant dans la suite 3, 5, 9, 15, 23~?
\end{question}

\begin{question}[$\star$ \iqroles{QT, Trader, Analyste} \iqfirm{teneur de marché, banque} \iqlevel{junior}]\label{iq:iv1:sequences-patterns-and-aptitude:3}
Quel est le nombre suivant dans la suite 1, 4, 27, 256~?
\end{question}

\begin{question}[$\star$ \iqroles{QT, Trader, Analyste} \iqfirm{teneur de marché, banque} \iqlevel{junior}]\label{iq:iv1:sequences-patterns-and-aptitude:4}
Quelle est la lettre suivante dans la suite O, T, T, F, F, S, S, E~?
\end{question}

\begin{question}[$\star$ \iqroles{QT, Trader, Analyste} \iqfirm{teneur de marché, banque} \iqlevel{junior}]\label{iq:iv1:sequences-patterns-and-aptitude:5}
Qu'est-ce qui vient ensuite dans la suite A1, C3, F6, J10, O15~?
\end{question}

\begin{question}[$\star$ \iqroles{QT, Trader, Analyste} \iqfirm{teneur de marché, banque} \iqlevel{junior}]\label{iq:iv1:sequences-patterns-and-aptitude:6}
Quel est le nombre suivant dans la suite 1, 20, 2, 17, 4, 14, 8, 11~?
\end{question}

\begin{question}[$\star$ \iqroles{QT, Trader, Analyste} \iqfirm{teneur de marché, banque} \iqlevel{junior}]\label{iq:iv1:sequences-patterns-and-aptitude:7}
Quel est le nombre suivant dans la suite 3, 6, 4, 8, 6, 12, 10~?
\end{question}

\begin{question}[$\star$ \iqroles{QT, Trader, Analyste, Risque} \iqfirm{teneur de marché, banque} \iqlevel{junior}]\label{iq:iv1:sequences-patterns-and-aptitude:8}
Une grille trois par trois contient neuf formes. La ligne 1 est composée d'un cercle blanc, d'un carré gris et d'un triangle noir. La ligne 2 est composée d'un triangle gris, d'un cercle noir et d'un carré blanc. La ligne 3 est composée d'un carré noir, d'un triangle blanc et d'une forme manquante. Quelle est la forme manquante~?
\end{question}

\begin{question}[$\star$ \iqroles{QT, Trader, Analyste, Risque, DS} \iqfirm{toutes} \iqlevel{junior}]\label{iq:iv1:sequences-patterns-and-aptitude:9}
Dans une société, tous les analystes ont un diplôme, et certaines personnes diplômées sont des traders. Peut-on en conclure que certains analystes sont des traders~? Répondez par oui ou par non, et justifiez votre réponse.
\end{question}

\begin{question}[$\star$ \iqroles{QT, Trader, Analyste, QD, Dev} \iqfirm{toutes} \iqlevel{junior}]\label{iq:iv1:sequences-patterns-and-aptitude:10}
Une liste contient cinq affirmations. La première affirme qu'exactement une affirmation de la liste est fausse, la deuxième qu'exactement deux sont fausses, la troisième exactement trois, la quatrième exactement quatre et la cinquième exactement cinq. Lesquelles sont vraies~?
\end{question}

\begin{question}[$\star$ \iqroles{QT, Trader, Analyste} \iqfirm{teneur de marché, banque} \iqlevel{junior}]\label{iq:iv1:sequences-patterns-and-aptitude:11}
Une règle transforme 1 en 2, 2 en 9, 3 en 28 et 4 en 65. En quel nombre transforme-t-elle 5~?
\end{question}

\begin{question}[$\star$ \iqroles{QT, Trader, Analyste} \iqfirm{teneur de marché, banque} \iqlevel{junior}]\label{iq:iv1:sequences-patterns-and-aptitude:12}
Dans un codage, le mot CAT s'écrit 24, et le mot DOG s'écrit 26. Comment s'écrit le mot FOX~?
\end{question}

\begin{question}[$\star$ \iqroles{QT, Trader, Analyste, Risque} \iqfirm{toutes} \iqlevel{junior}]\label{iq:iv1:sequences-patterns-and-aptitude:13}
Une voiture parcourt 120~km à 40~km/h, puis revient aussitôt à 60~km/h. Quelle est sa vitesse moyenne sur l'ensemble du trajet~?
\end{question}

\begin{question}[$\star$ \iqroles{QT, Trader, Analyste, Risque} \iqfirm{toutes} \iqlevel{junior}]\label{iq:iv1:sequences-patterns-and-aptitude:14}
Le tuyau A remplit seul un réservoir en 6~heures, et le tuyau B le remplit seul en 3~heures. Une vidange vide seule un réservoir plein en 4~heures. Le réservoir est vide, et les deux tuyaux ainsi que la vidange sont ouverts ensemble. Combien de temps le réservoir met-il pour se remplir~?
\end{question}

\begin{question}[$\star$ \iqroles{QT, Trader, Analyste, Risque, DS} \iqfirm{toutes} \iqlevel{junior}]\label{iq:iv1:sequences-patterns-and-aptitude:15}
Dans une promotion de 40 étudiants, 25 suivent des mathématiques, 18 suivent de la physique et 7 ne suivent ni l'un ni l'autre. Combien suivent les deux~? Si l'on ne vous disait pas combien suivent aucune de ces matières, quel intervalle de valeurs serait possible~?
\end{question}

\begin{question}[$\star$ \iqroles{QT, Trader, Analyste, Risque} \iqfirm{banque, hedge fund} \iqlevel{junior}]\label{iq:iv1:sequences-patterns-and-aptitude:16}
Trois produits présentent les revenus et les coûts suivants sur deux trimestres. La marge désigne le revenu moins le coût, exprimé en pourcentage du revenu. Quel produit a vu sa marge progresser le plus en points de pourcentage, et quel produit a vu son profit progresser le plus en dollars~?

\smallskip
\begin{tabular}{lrrrr}
\hline
Produit & Revenu T1 & Coût T1 & Revenu T2 & Coût T2 \\
\hline
X & 200 & 150 & 250 & 180 \\
Y & 400 & 320 & 500 & 395 \\
Z & 120 & 96 & 150 & 111 \\
\hline
\end{tabular}
\end{question}

\begin{question}[$\star\star$ \iqroles{QT, Trader, Analyste} \iqfirm{teneur de marché, banque} \iqlevel{junior}]\label{iq:iv1:sequences-patterns-and-aptitude:17}
Quel est le nombre suivant dans la suite 6, 15, 35, 77, 143~?
\end{question}

\begin{question}[$\star\star$ \iqroles{QT, Trader, Analyste} \iqfirm{teneur de marché, banque} \iqlevel{junior}]\label{iq:iv1:sequences-patterns-and-aptitude:18}
Quel est le nombre suivant dans la suite 19, 29, 40, 44, 52, 59~?
\end{question}

\begin{question}[$\star\star$ \iqroles{QT, Trader, Analyste} \iqfirm{teneur de marché, banque} \iqlevel{junior}]\label{iq:iv1:sequences-patterns-and-aptitude:19}
Quel est le nombre suivant dans la suite 2, 3, 5, 4, 9, 7, 8, 27, 9, 16~?
\end{question}

\begin{question}[$\star\star$ \iqroles{QT, Trader, QD, Dev} \iqfirm{teneur de marché, banque} \iqlevel{junior}]\label{iq:iv1:sequences-patterns-and-aptitude:20}
Quel est le nombre suivant dans la suite 1, 2, 3, 4, 5, 10, 11, 12, 13, 14, 15, 20~?
\end{question}

\begin{question}[$\star\star$ \iqroles{QT, Trader, QD, Dev, Analyste} \iqfirm{teneur de marché, banque} \iqlevel{junior}]\label{iq:iv1:sequences-patterns-and-aptitude:21}
Chaque case d'une grille trois par trois représente un carré dont seuls certains des quatre côtés sont tracés. Ligne 1~: la première case a les côtés haut et gauche, la deuxième a les côtés gauche et droit, la troisième a les côtés haut et droit. Ligne 2~: la première a les côtés bas et droit, la deuxième les côtés haut et droit, la troisième les côtés haut et bas. Ligne 3~: la première a les côtés haut, gauche et bas, la deuxième les côtés gauche et droit, et la troisième est manquante. Quels côtés possède la case manquante~?
\end{question}

\begin{question}[$\star\star$ \iqroles{QT, Trader, Analyste} \iqfirm{teneur de marché, banque} \iqlevel{junior}]\label{iq:iv1:sequences-patterns-and-aptitude:22}
Chaque ligne d'un tableau contient trois nombres, et le troisième se déduit des deux premiers par la même règle. Les lignes sont (2, 3, 13), (4, 1, 17), (1, 5, 26) et (3, 3, 18). Quel nombre complète la ligne (5, 2, ~?)
\end{question}

\begin{question}[$\star\star$ \iqroles{QT, Trader, QD, Dev} \iqfirm{toutes} \iqlevel{junior}]\label{iq:iv1:sequences-patterns-and-aptitude:23}
Chacune de trois personnes dit toujours la vérité ou ment toujours. A affirme que B est un menteur. B affirme que C est un menteur. C affirme que A et B sont tous deux des menteurs. Qui dit la vérité~?
\end{question}

\begin{question}[$\star\star$ \iqroles{QT, Trader, Analyste, QD, Dev} \iqfirm{toutes} \iqlevel{junior}]\label{iq:iv1:sequences-patterns-and-aptitude:24}
Cinq coureurs terminent une course sans égalité. C termine devant A, et A termine devant B. E termine juste devant D. B ne termine pas dernier. Quel est l'ordre d'arrivée~?
\end{question}

\begin{question}[$\star\star$ \iqroles{QT, Trader, Analyste} \iqfirm{toutes} \iqlevel{junior}]\label{iq:iv1:sequences-patterns-and-aptitude:25}
Anna a deux fois l'âge qu'avait Ben lorsque Anna avait l'âge actuel de Ben. Leurs âges actuels totalisent 63. Quel âge a chacun d'eux~?
\end{question}

\begin{question}[$\star\star$ \iqroles{QR, QT, DS, Analyste} \iqfirm{toutes} \iqlevel{junior}]\label{iq:iv1:sequences-patterns-and-aptitude:26}
Une question de test affiche 1, 2, 4, 8, 16 et demande le nombre suivant. Un candidat répond 31 et peut présenter une formule qui s'accorde avec les cinq termes et qui donne 31. La question est-elle défectueuse~? Quelle réponse donneriez-vous à un test, et pourquoi~?
\end{question}

\begin{question}[$\star\star\star$ \iqroles{QT, Trader, Analyste} \iqfirm{teneur de marché, banque} \iqlevel{junior}]\label{iq:iv1:sequences-patterns-and-aptitude:27}
Quel est le nombre suivant dans la suite 1, 4, 9, 7, 7, 9, 13, 10~?
\end{question}

\begin{question}[$\star\star\star$ \iqroles{QR, QT, Trader} \iqfirm{teneur de marché, hedge fund} \iqlevel{junior}]\label{iq:iv1:sequences-patterns-and-aptitude:28}
Quel est le nombre suivant dans la suite 1, 2, 3, 7, 22, 155~?
\end{question}

\begin{question}[$\star\star\star$ \iqroles{QR, QT, Trader, QD} \iqfirm{teneur de marché, hedge fund} \iqlevel{junior, confirmé}]\label{iq:iv1:sequences-patterns-and-aptitude:29}
Quel est le nombre suivant dans la suite 1, 1, 2, 5, 14, 42~? Donnez la règle. Deux autres suites bien connues commencent de la même façon~: pourquoi ne conviennent-elles pas~?
\end{question}

\begin{question}[$\star\star\star$ \iqroles{QR, QT, Trader, QD, Dev} \iqfirm{teneur de marché, hedge fund} \iqlevel{junior, confirmé}]\label{iq:iv1:sequences-patterns-and-aptitude:30}
Quel est le nombre suivant dans la suite 1, 2, 4, 8, 16, 22, 26, 38, 62, 74~?
\end{question}
